\subsection{Unified Reference Protocol and Completed Follow-up}
\label{sec:appendix:unified_v2}

The unified v2 comparison replaces the historical reference-recipe leaderboard
in Table~\ref{tab:main}, not the separate prior-controlled and paired-RNA
diagnostics. The protocol was fixed before generating the v2 outer predictions;
the existing outer folds had already been examined, so this is an exploratory
follow-up rather than an untouched confirmatory evaluation.

\paragraph{Common data and isolation.}
The input is the barcode-aligned raw APT count matrix with 293 columns in
fixed source order. Each count is transformed by log1p, then standardized
feature-wise using only the current fit partition. All annotated benchmark
cells are used, without patient caps, class weights, patient weights or
replacement sampling. Equal patient weight is used for scoring, not training.
For test fold $f$, validation is $(f+1)\bmod 5$ and the other three folds
form the inner fit. A new scaler and a newly initialized model are fitted
on all development patients after selection. Outer-test cells are never used
for fitting the scaler, choosing a configuration or determining stopping.
Table~\ref{tab:unified_fold_counts} gives exact observed counts.
\begin{table}[H]
\centering\small
\setlength{\tabcolsep}{4pt}
\begin{tabular}{ccccc}
\toprule
Fold & Inner fit & Validation & Outer refit & Test \\
\midrule
1 & 24 / 226,431 & 8 / 68,381 & 32 / 294,812 & 8 / 66,980 \\
2 & 24 / 207,172 & 8 / 86,239 & 32 / 293,411 & 8 / 68,381 \\
3 & 24 / 206,704 & 8 / 68,849 & 32 / 275,553 & 8 / 86,239 \\
4 & 24 / 221,600 & 8 / 71,343 & 32 / 292,943 & 8 / 68,849 \\
5 & 24 / 223,469 & 8 / 66,980 & 32 / 290,449 & 8 / 71,343 \\
\bottomrule
\end{tabular}
\caption{Exact unified v2 patient / cell counts, reconstructed from input IDs and frozen partitions and checked against every recipe. All methods share these counts. Validation is the next cyclic outer fold; remaining patients form the inner fit. Outer refit combines inner fit and validation. Counts include all annotated benchmark cells, without a cap.}
\label{tab:unified_fold_counts}
\end{table}


\paragraph{Model-specific search spaces.}
LR uses unweighted multinomial logistic regression with LBFGS, a maximum of
1{,}000 iterations and $C\in\{0.01,0.1,1,10\}$. XGBoost uses CPU histogram
trees (four requested threads), depth $\{3,6\}$, learning rate
$\{0.03,0.1\}$, and tree counts $\{100,300,500\}$, with row and column
subsampling both 0.8. All twelve combinations are fitted separately; this
is not continuous tree-wise early stopping. LR and XGBoost independently
select coarse and fine configurations by fixed-ontology validation SB-F1,
then refit using the chosen $C$ or tree configuration.

The plain MLP has 256/128-unit ReLU hidden layers, dropout 0.1 and separate
linear fine/coarse heads. All Transformer variants share the existing APT
feature-identity/value tokenizer and pre-norm encoder, with hidden dimension
128, two layers, four attention heads, feed-forward dimension 256 and dropout
0.1. These dimensions differ from the historical 256-dimensional,
four-layer references in Appendix~\ref{sec:method}.
Neural models use AdamW with the full grid
$\mathrm{lr}\in\{3\times10^{-4},10^{-3}\}$ and
$\mathrm{weight\ decay}\in\{10^{-5},10^{-3}\}$. Batch size is 1{,}024 for
MLP and 128 for Transformers. Each candidate trains for at most 50 epochs,
with validation fine SB-F1 after every epoch and patience eight. Strict
improvement updates the best epoch; ties retain the earliest epoch and first
candidate. Fresh outer refits run for that selected epoch count without test
monitoring. The v2 loop uses float32, no gradient clipping and no learning-rate
schedule. The historical mixed-precision and contrastive warm-up settings
are not used here.

\paragraph{Objectives and component scope.}
MLP, Flat and Cascade minimize unweighted fine CE plus 0.5 coarse CE.
Cascade retains the global path, lineage experts and log compatibility term
in Eq.~\ref{eq:cascade}, with temperature 1, expert scale 1, $\epsilon=0.1$
and constant $\gamma=0.5$ (no ramp). Flat bypasses experts and routing.
The wrapper retains the common projection head and unused expert parameters
for Flat; effective capacity and compute are therefore not identical.
HCE instead minimizes the leaf and parent negative log subtree masses, with
weights 1 and 0.5, over the joint 32-node softmax. Fine prediction selects
among leaves; parent prediction sums each internal node and its descendants.
No main-table model uses disease labels as training supervision; the inherited
outer folds remain disease-stratified.

For the additional Flat/Cascade contrastive variants, normalize projected
embeddings to unit length. With $P_i$ the same-subtype, different-disease
training examples in the mini-batch, the added objective is
\begin{equation}
0.1\mathcal L_{\rm con},\qquad
\mathcal L_{\rm con}=-\frac{1}{|V|}\sum_{i\in V}\frac{1}{|P_i|}
\sum_{j\in P_i}\log\frac{\exp(u_i^\top u_j/0.1)}
{\sum_{a\ne i}\exp(u_i^\top u_a/0.1)}.
\end{equation}
Here $V$ contains anchors with positives; a mini-batch without valid anchors
contributes zero. All non-self examples enter the denominator, including
same-subtype/same-disease examples. This common v2 loss is distinct from the
historical positive-mass objective in Appendix~\ref{sec:method:contrastive}.
No disease labels are required at inference. Each of the four recipes selects
its own validation optimum under the same grid: the resulting contrasts
compare whole training/selection procedures, not a routing-only intervention.

\paragraph{Seeds, scoring and intervals.}
Selection is performed once per model/fold (and per classical task) with
seed 17; randomized final fits use seeds 17, 29 and 43. LR is deterministic
and is fitted once per task/fold. A seed controls model initialization and,
for neural models, a without-replacement cell permutation each epoch; it
does not change the folds or number of unique training cells. Scores first
pool patient-normalized confusion matrices over the complete OOF set, then
average the resulting seed-specific F1 values, not predictions.
Fine and coarse ontologies always contain 27 and five labels respectively,
with zero-denominator F1 set to zero.

The bootstrap uses 5{,}000 replicates and RNG seed 20260915. Each replicate
draws 40 patients with replacement and three seed indices with replacement.
The same patient draw and matched seed slots are used across all conditions;
F1 is computed separately in each seed slot and then averaged. LR reuses its
one fit, so its uncertainty reflects only patient resampling. Paired
differences and the $2\times2$ interaction are computed within each replicate,
with percentile 95\% intervals. These are pointwise, exploratory intervals,
conditional on the chosen configurations, folds and finite seed set; they
do not repeat model selection or correct for multiplicity. Shared development
sets also prevent interpreting the OOF patients as independent fitted models.
No equivalence margin or equivalence test is supplied, so overlapping or
zero-containing intervals do not establish equivalence.

\paragraph{Deployable versus privileged decoding.}
Let $p_k$ be the frozen fine probabilities, $q_m$ the frozen parent
probabilities and $g(k)$ the fixed parent mapping. D0 takes $\arg\max_k p_k$.
D1 restricts that argmax to $g(k)=\arg\max_m q_m$.
D2 takes the argmax of
$q_{g(k)}p_k/\sum_{j:g(j)=g(k)}p_j$ (a $10^{-30}$ denominator floor).
D4 restricts the argmax to the true reference parent. No weights change,
and D4 is not used to select a model or report deployable performance.
These rules use subtype probabilities and a fixed ontology; they are not
new fitted hierarchical models. All conditions, including unfavorable
decoding results, are retained in Table~\ref{tab:unified_decoding}.
\begin{table}[H]
\centering\small
\setlength{\tabcolsep}{4pt}
\begin{tabular}{lcccc}
\toprule
Model & Ordinary D0 & Hard D1 & Consistent D2 & Oracle D4 \\
\midrule
Logistic Regression & 0.134 & 0.110 & 0.116 & 0.360 \\
XGBoost & 0.134 & 0.126 & 0.131 & 0.362 \\
Plain MLP & 0.137 & 0.131 & 0.134 & 0.376 \\
Flat Transformer & 0.147 & 0.137 & 0.142 & 0.371 \\
HCE Transformer & 0.146 & 0.142 & 0.146 & 0.371 \\
Soft-cascade (no disease) & 0.148 & 0.136 & 0.143 & 0.371 \\
Flat + contrastive & 0.148 & 0.137 & 0.143 & 0.372 \\
Cascade + contrastive & 0.150 & 0.138 & 0.143 & 0.376 \\
\bottomrule
\end{tabular}
\caption{Frozen-score decoding diagnostics for all unified v2 recipes (fine SB-F1, seed means). D1 restricts candidates to the predicted parent; D2 combines predicted parent probabilities with within-parent subtype probabilities. Both are deployable post-processing rules using the same frozen scores. D4 uses the true parent and is privileged, not a deployable score or an information upper bound. LR/XGBoost use separately selected parent models; neural references use their auxiliary parent outputs, with subtree probabilities for HCE. No retraining or test-based tuning of the decoding rules is performed. Full scores, seed values and paired intervals are released in the v2 CSV summaries.}
\label{tab:unified_decoding}
\end{table}


\paragraph{Artifact validation and HCE readout correction.}
All 130 final score arrays are checked against the original cell IDs,
patient IDs, labels, fixed class order and outer-fold membership without
intersection filtering. Recorded fit counts and input hashes agree across
recipes; inner and outer standardization statistics are independently
reconstructed from the raw inputs. The selected candidate and earliest best
epoch are checked against saved validation histories. During aggregation,
the initial HCE export was found to normalize internal-node logits alone
for its parent readout. This does not match its subtree-supervised objective.
Fifteen frozen HCE checkpoints were replayed to recover joint-node subtree
probabilities, without retraining, reselection or changing fine predictions.
Only coarse, D1 and D2 results are affected; D0 and D4 remain unchanged.
The original exports, checkpoint hashes and correction record are retained.
This was a readout correction, not a new training experiment.

\paragraph{Compute and reproducibility artifacts.}
Runs were split between vllab11 and vllab15, both using RTX 6000 Ada GPUs for
neural fits; LR/XGBoost remained on CPU. Per-fit audit records identify
hardware, exact counts, seeds, selected parameters and search/refit timing.
Table~\ref{tab:unified_compute} summarizes realized budgets.
The repository directory \texttt{analysis/unified\_v2/} contains the frozen
protocol, training and scoring scripts, source-validation script, HCE replay,
machine-readable stage inventory, per-seed scores, all paired contrasts and
source hashes. These supplement rather than resolve the existing data-access
requirements. Historical scores and diagnostics remain separately identified
in the following appendices; they are not pooled with v2 fits.
\begin{table}[H]
\centering\small
\setlength{\tabcolsep}{4pt}
\begin{tabular}{lcccc}
\toprule
Recipe & Candidate fits & Refit epochs & Search h & Refit h \\
\midrule
Logistic Regression & 40 & --- & 0.53 & 0.20 \\
XGBoost & 120 & --- & 4.32 & 2.31 \\
Plain MLP & 20 & 15--26 & 0.08 & 0.04 \\
Flat Transformer & 20 & 14--45 & 4.73 & 3.79 \\
HCE Transformer & 20 & 14--37 & 4.10 & 3.34 \\
Soft-cascade (no disease) & 20 & 11--39 & 5.07 & 3.84 \\
Flat + contrastive & 20 & 11--47 & 5.85 & 6.19 \\
Cascade + contrastive & 20 & 19--46 & 6.23 & 7.03 \\
\bottomrule
\end{tabular}
\caption{Realized v2 search/refit records summed across folds (and across tasks for LR/XGBoost). Neural search uses four candidates per fold; LR uses four and XGBoost twelve per task/fold. Epoch ranges are selected stopping epochs, not the epochs actually consumed by all search trials. Times are accumulated timed-block wall hours, not elapsed project time or GPU-hours: parallel jobs overlap, hosts differ, and input/scaler preparation is outside these blocks. Neural runs used RTX 6000 Ada GPUs; classical runs used CPU. This is a reproducibility record, not a cross-hardware efficiency ranking. HCE correction inference is excluded.}
\label{tab:unified_compute}
\end{table}

